\subsection{Backward verifier}\label{sec:impl-backward}
The backward verifier
(\texttt{src/v5/backward\_shape.py},
\texttt{src/v5/grad\_flag\_verifier.py})
consumes a forward graph annotated with the refinement types of
\Cref{sec:calculus} and certifies, without execution, that
(i)~\texttt{loss.backward()} would produce gradients of the right shape
modulo \texttt{sum\_to(grad, primal.shape)} reduction, and (ii)~the set
$H$ of parameters reverse-reachable from \texttt{loss} along
non-detached, non-\texttt{no\_grad}-context, \texttt{requires\_grad=True}
edges equals the user-declared expected-to-learn set $E$. The check is
a reverse-topological traversal that mirrors PyTorch's per-Function
shape contract \verb|grad_outputs[i].shape == y_i.shape|. The
\emph{grad-flag} component additionally detects three canonical
silent-zero-grad bug classes: in-place writes on a leaf with
\texttt{requires\_grad=True} (B3), no leaf in the trace has
\texttt{requires\_grad=True} (B4), and parameters used only inside a
\texttt{with torch.no\_grad():} block or after \texttt{.detach()} (B1).
These are reported as \emph{Refuted} verdicts with witness traces.

\paragraph{Empirical agreement.}
On a 500-model property test (random small \texttt{nn.Module}s drawn
from a grammar), the static verifier agrees with PyTorch's runtime
behaviour on $500/500$ models (gradient presence and gradient shape
both match). On 8 hand-curated case studies modelled on canonical
PyTorch issues (no-grad scope, \texttt{.detach()},
\texttt{requires\_grad=False}, in-place-on-leaf, missing leaf, in-place
alias mutating a saved tensor, non-scalar loss, unused parameter
silently skipped by \texttt{optimizer.step()}), the verifier correctly
catches $8/8$. On a 50-model false-positive sweep on clean models, the
verifier reports $0$ spurious refutations. All 20 unit tests
pass. Numbers are in \texttt{experiments\_v5/track\_D\_results.json}.

